\section{Introduction}\label{sec:introduction}

The minimal model program seeks birational models on which the canonical
divisor has a definite sign: nef on a minimal model, or negative along the
fibres of a Mori fibre space. For a generalized pair, the adjoint divisor
also contains the trace of nef data on a higher birational model. This
generalized-pair formalism of Birkar--Zhang
\cite[Definition~1.4]{BirkarZhang} accommodates the moduli terms arising
in the canonical bundle formula \cite{AmbroCBF}. The existence question
therefore asks for a birational model while keeping those nef data fixed.

We prove the following form of the existence conjecture. Our discrepancy
convention is the log convention: the discrepancy of a component of a
crepant boundary of coefficient $b$ is $1-b$.

\begin{theorem}\label{thm:main}
Let $k$ be an algebraically closed field of characteristic zero, and let
$(X,B+M)$ be a projective generalized log canonical $\Q$-pair. Thus $X$ is
normal and integral, $B$ is an effective rational boundary, and $M$ is the
trace of a nef rational Cartier divisor on a higher projective birational
model, with $K_X+B+M$ $\Q$-Cartier.

There are a normal integral projective variety $Y$ and a birational map
$\phi\colon X\dashrightarrow Y$ whose inverse contracts no divisors such
that, writing $B_Y=\phi_*B$ and taking $M_Y$ to be the trace of the same
nef b-divisor, the generalized pair $(Y,B_Y+M_Y)$ is log canonical and
\[
 A_{X,B+M}(E)\le A_{Y,B_Y+M_Y}(E)
\]
for every prime divisor $E$ over their common function field. The
inequality is strict for every divisor on $X$ contracted by $\phi$.
More precisely, the following alternatives are determined by $D=K_X+B+M$:
\begin{enumerate}
\item if $D$ is pseudo-effective, then $K_Y+B_Y+M_Y$ is nef;
\item if $D$ is not pseudo-effective, there is a projective contraction $h\colon Y\to T$ to a normal
projective variety with connected fibres such that
\[
 \dim T<\dim Y,\qquad \rho(Y/T)=1,
 \qquad -(K_Y+B_Y+M_Y)\text{ is ample over }T.
\]
\end{enumerate}
\end{theorem}

No $\Q$-factoriality is imposed on $X$ or $Y$. The discrepancy comparison
uses the original nef b-divisor and the stated log normalization.
The assertion is the existence form of the minimal model conjecture
\cite[Conjecture~1.1]{GongyoMinimalModel}, with fixed rational generalized
nef data. It gives a terminating choice of birational program; it makes
no assertion that every arbitrary MMP terminates. The nef conclusion of
Theorem~\ref{thm:main} does not itself include semiampleness or the
existence of a nonzero pluricanonical section.

\begin{corollary}[Ordinary log canonical pairs]\label{cor:ordinary}
Let $(X,B)$ be a projective log canonical $\Q$-pair over an algebraically
closed field of characteristic zero. There is a birational map
$\phi:X\dashrightarrow Y$ extracting no divisors, with $B_Y=\phi_*B$,
such that $(Y,B_Y)$ is log canonical and
\[
 A_{X,B}(E)\leq A_{Y,B_Y}(E)
\]
for every prime divisor over the common function field, strictly for
prime divisors on $X$ contracted by $\phi$. If $K_X+B$ is
pseudo-effective, $K_Y+B_Y$ is nef; otherwise $Y$ admits a Mori fibre
contraction with relative Picard number one and relatively ample
$-(K_Y+B_Y)$.
\end{corollary}
\begin{proof}
Apply Theorem~\ref{thm:main} to the fixed zero nef b-divisor.
\end{proof}

In the pseudo-effective ordinary case, the companion log abundance theorem
\cite[Theorem~1.1]{OpenAILogAbundance2026} makes $K_Y+B_Y$ semiample on
the same endpoint: $(Y,B_Y)$ is a projective lc pair with effective rational
boundary and nef $\Q$-Cartier adjoint. The proof of Corollary~12.1(i) in the same
reference already uses this application of Corollary~\ref{cor:ordinary}.
This separate consequence gives no semiampleness assertion for nonzero
generalized nef data.

The ordinary case illustrates the role of the generalized data. When
$M=0$, the crepant equation defining discrepancies contains only the
canonical divisor and the boundary. For a generalized pair one pulls
back the same nef divisor from its determination on every common
resolution; it cancels in comparisons between successive models.
The local phase theorem concerns ordinary Gorenstein klt germs in both
cases. The generalized structure enters the transfer from the smooth
relative problem to the final pair, where the NQC MMP applies because
the fixed nef data are rational.

\subsection{Background and significance}

The extremal-ray and contraction approach to the minimal model problem
originates in Mori's work \cite{Mori1982}; the subsequent theory of
singular pairs is developed systematically in Koll\'ar--Mori
\cite{KM98}. Birkar--Cascini--Hacon--McKernan \cite{BCHM}
established existence and finite generation for klt pairs in the big
setting. Their results supply the flips, positive-boundary programs and
finiteness of ample models used in this paper. The remaining
pseudo-effective case requires an argument that does not assume
bigness of the adjoint or boundary.

Weak Zariski decompositions organize one route from that problem to
existence. Birkar \cite{BirkarWZD} related weak decompositions and log
minimal models under lower-dimensional MMP hypotheses. Han--Li
\cite{HanLiWZD} developed the generalized-pair version, and
Lazi\'c--Tsakanikas \cite{LazicTsakanikasMM} reduced ordinary-lc
existence to the relative smooth problem. Their later special-MMP
results \cite{LazicTsakanikasSpecial} give a $\Q$-factorial NQC
precursor to the transfer used here. Tsakanikas--Xie
\cite{TsakanikasXie} remove that ambient factoriality restriction
using the generalized MMP foundations, including the flip results of
Hacon--Liu and Liu--Xie and the contractions of Xie
\cite{HaconLiu,LiuXie,XieContractions}. Their terminating programs have
explicit existence or non-pseudo-effectivity hypotheses. We establish
the smooth input and its relative weak Zariski decompositions before
applying those theorems; the transfer itself is credited to them.
Hu--Liu subsequently established flips without the NQC hypothesis
\cite[Theorems~1.1--1.2]{HuLiuNonNQC}; the terminating transfer used here
remains in the NQC setting.

The local argument uses the normalized-volume theory initiated by Li
\cite{LiMinimizing}. Blum proved existence of minimizing valuations
\cite[Main Theorem and Section~7]{BlumExistence}, Xu proved their
quasi-monomiality \cite[Theorem~1.2]{XuQuasiMonomial}, and Xu--Zhuang
established uniqueness and finite generation
\cite{MinimizerUniqueness,StableDegeneration}. These results supply
the invariant minimizing valuation and its finitely generated
klt cone. Section~\ref{sec:cones} establishes the additional equivariant
rational regrading and Gorenstein Rees properties needed here.
Birkar's complements, boundedness of Fano varieties and effective
birationality for polarized varieties
\cite{BirkarComplements,BAB,BirkarPolarised} provide the projective
numerical inputs. The new product estimate links those inputs to
controlled valuation selection and exact character congruences.

The relative estimates build on Ambro's canonical bundle formula
\cite{AmbroCBF} and the b-semiampleness theorem of
Bakker--Filipazzi--Mauri--Tsimerman \cite{BFMT}. Their bounded-family
effectiveness results give a related precedent under additional
hypotheses, including klt singularities
\cite[Corollaries~7.8--7.9]{BFMT}; the reference form here may be lc.
Each use of the canonical bundle formula checks generic effectivity,
discrepancy rank one and the adjoint pullback identity.
Uniformity then comes from the finitely many prepared families, not
from an unstated uniform index in the b-semiampleness theorem.
Toroidal preparation
uses Abramovich--Denef--Karu and Abramovich--Karu
\cite{ADK,AbramovichKaruRefinement}; the fibre comparisons also use
Koll\'ar's linking of lc centres \cite{KollarSources} and Birkar's
very-exceptional-divisor method \cite{BirkarFlips}. Valuative
retraction follows the framework of Jonsson--Musta\c t\u a
\cite{JonssonMustata}, while logarithmic Cartier and
Deligne--Illusie theory supply the positive-characteristic selection
step \cite{Cartier1957,Katz1970,DeligneIllusie}. The global use of
asymptotic multiplier ideals comes from Demailly--Ein--Lazarsfeld
\cite{DEL}. These established methods are distinguished below from
the uniform lifting, mixing, rounding and termination arguments
proved in this paper.

\subsection{The local input and the global contradiction}

The central local statement is Theorem~\ref{thm:phase}. For a sequence of
Gorenstein klt germs of fixed dimension with cyclic symmetry, it selects
divisorial valuations of bounded discrepancy whose values recover, up
to one fixed positive integer, the character angles of a chosen group
element. The congruence holds for every rational semi-invariant, and
the restricted valuation on the invariant field is integral-valued.

To apply it, suppose that a canonical MMP with scaling, starting from a
smooth projective variety, is infinite. Uniformly generated asymptotic
multiplier ideals bound the accumulated discrepancy changes. On
integral-valued divisorial tests of bounded original discrepancy, the
set of finite limits of these changes is countable. The phase theorem
on cyclic index-one covers then rules out unbounded canonical indices:
otherwise arbitrary character angles would lie in a countable set.
With one common Cartier multiple fixed, higher-direct-image vanishing
on late flips and effective base point freeness force a divisor to be
both globally generated and negative on a contracted curve. This
contradiction proves the smooth absolute input. The relative reduction
and descent to the stated field are included in
Section~\ref{sec:global}.

\subsection{Uniform estimates and the structure of the proof}

The exposition begins with the global deduction, then develops its local
input, and finally proves the deeper uniform estimates used there.
There are three useful reading routes.

For the global conclusion, Theorem~\ref{thm:phase} implies smooth
absolute termination (Proposition~\ref{prop:smooth-termination}).
Proposition~\ref{prop:relative-wzd} supplies the relative smooth
weak Zariski decomposition, after which the generalized existence
theorems and Proposition~\ref{prop:field-descent} give
Theorem~\ref{thm:main} over the stated field.

For the phase construction, the cone estimates of
Proposition~\ref{prop:cone-estimates} and the degree-product estimate
of Proposition~\ref{prop:product} lead to constrained valuation
optimization, logarithmic Cartier selection and exact rounding in
Sections~\ref{sec:valuative}--\ref{sec:rounding}. The characteristic
reduction occurs only after the germ and its required finite data have
been fixed. A finite recursion chooses all bounds before the final
integer is chosen, and the resulting congruence holds for every
semi-invariant. Integral-valued divisorial valuations in this argument
need not be primitive.

For the uniform estimates, bounded presentations, actual trait
calculations and lc-centre incidence in
Sections~\ref{sec:bounded-presentations}--\ref{sec:incidence} prove
Theorem~\ref{thm:bounded-fibre}. Only its reference form has bounded
index. Tower reduction then gives Theorem~\ref{thm:lifting-mixing},
which enters the integral jump estimate and hence the product bound.
Ordinary fibre-component labels are retained along with cone weights:
weights alone do not identify the required actual valuation. Saturated
face lattices eliminate hidden orbit covers. The effective dlt
construction in Lemma~\ref{inc:effective-model} supplies the lc-incidence
input by eliminating a very exceptional divisor on a fixed finite
model, independently of Theorem~\ref{thm:main}.

The following locators make the deferred proofs explicit.
\begin{center}
\small
\begin{tabular}{lll}
\toprule
Result & Statement & Proof \\
\midrule
Phase theorem & Section~\ref{sec:global} & Sections~\ref{sec:cones}--\ref{sec:rounding} \\
Integral jump & Section~\ref{sec:product} & Section~\ref{sec:integral-jump} \\
Lifting and Mixing & Section~\ref{sec:integral-jump} & Section~\ref{sec:tower} \\
Bounded fibre & Section~\ref{sec:tower} & Sections~\ref{sec:bounded-presentations}--\ref{sec:incidence} \\
\bottomrule
\end{tabular}
\end{center}
The arrows in Figure~\ref{fig:dependencies} give the logical order of
these proofs, including their later-section inputs.

\begin{figure}[p]
\centering
\begin{tikzpicture}[
 box/.style={draw,rounded corners=2pt,align=center,text width=5.8cm,
 minimum height=0.8cm,inner sep=5pt,font=\small},
 side/.style={box,text width=3.7cm},
 >={Stealth[length=5pt]},node distance=0.5cm and 0.6cm]
 \node[box] (prepare) {Bounded presentations, traits and lc incidence\\Sections~\ref{sec:bounded-presentations}--\ref{sec:incidence}};
 \node[box,below=of prepare] (bounded) {Bounded-fibre estimates\\Theorem~\ref{thm:bounded-fibre}};
 \node[box,below=of bounded] (tower) {Tower reduction and Lifting/Mixing\\Section~\ref{sec:tower}; Theorem~\ref{thm:lifting-mixing}};
 \node[box,below=of tower] (jump) {Integral jump estimate\\Proposition~\ref{prop:integral-jump}};
 \node[box,below=of jump] (product) {Product of degree scales\\Proposition~\ref{prop:product}};
 \node[side,right=of product] (cone) {Equivariant cone estimates\\Proposition~\ref{prop:cone-estimates}};
 \node[box,below=of product] (select) {Valuation optimization\\$\to$ Cartier selection\\$\to$ exact rounding\\Sections~\ref{sec:valuative}--\ref{sec:rounding}};
 \node[box,below=of select] (phase) {Local phase theorem\\Theorem~\ref{thm:phase}};
 \node[box,below=of phase] (global) {Smooth absolute $\to$ smooth relative\\$\to$ generalized existence\\$\to$ field descent\\Section~\ref{sec:global}; Theorem~\ref{thm:main}};
 \foreach \a/\b in {prepare/bounded,bounded/tower,tower/jump,jump/product,product/select,select/phase,phase/global}
   \draw[->] (\a)--(\b);
 \draw[->] (cone)--(product);
 \draw[->] (cone.south) |- (select.east);
\end{tikzpicture}
\caption{Logical dependencies. The preparation group includes the trait
minimum identity used for the bounded base form; no bounded moduli index
is assumed in proving that identity. The tower box also includes the
geometric step used directly in the integral jump argument.}
\label{fig:dependencies}
\end{figure}
